%% kapa: exam=2 round=26 year=2015 period=3 subject=compensation_law subject_ko=감정평가 및 보상법규 seq=3 number=3 points=20
%% source: Q-Net 감정평가사 기출문제 (cst00309) artlSeq=5149601 file=extracted/1008/5149601/2015년 제26회 감정평가사 2차시험 기출문제/3교시(감정평가 및 보상법규).pdf
%% conversion: kordoc 4.24.51 -> tools/gen2.py (auto); review=proofread-gpt-sol; flags=none
\begin{longquestion}{【문제 3】}
甲은 C시 소재 전(田) 700 m\textsuperscript{2}(이하 ‘이 사건 토지’라고 한다)의 소유자로서, 여관 신축을 위하여 부지를 조성하였는데, 진입로 개설비용 3억 원, 옹벽 공사비용 9천만 원, 토목설계비용 2천만 원, 토지형질변경비용 1천만 원을 각 지출하였다. 그런데 건축허가를 받기 전에 국토교통부장관이 시행하는 고속도로건설공사에 대한 사업인정이 2014년 7월 15일 고시되어 이 사건 토지 중 500 m\textsuperscript{2}(이하 ‘이 사건 수용 대상 토지’라고 한다)가 공익사업시행지구에 편입되었고, 2015년 7월 17일 관할 토지수용위원회에서 수용재결이 있었다. 그 결과 이 사건 토지에서 이 사건 수용 대상 토지를 제외한 나머지 200 m\textsuperscript{2}(이하 ‘이 사건 나머지 토지’라고 한다)는 더 이상 여관 신축의 용도로는 사용할 수 없게 되어 그 부지조성 비용은 이 사건 나머지 토지의 정상적인 용도에 비추어 보았을 때에는 쓸모없는 지출이 되고 말았다. 이에 甲은 이 사건 나머지 토지에 들인 부지조성 비용에 관하여 손실보상의 지급을 청구하고자 한다. 다음 물음에 답하시오. (20점) \par
1) 위 청구권의 법적 근거에 관하여 설명하시오. (10점) \par
2) 甲은 다른 절차를 거치지 않고 바로 국가를 상대로 손실보상을 청구하는 소송을 제기할 수 있는가? (10점) \par
\end{longquestion}
